% Irodov Part III §3.4 current 161--175; user PDF pp.123--125,
% printed pp.114--116. Source 168 incomplete English clause, noted.
\electrocurrentproblem{161}{114}
两电极形状任意，浸没在均匀、弱导电介质中，
介质电阻率 $\rho$、相对介电常数 $\varepsilon$。
求两电极间漏电电阻 $R$ 与互电容 $C$ 的乘积。
\exsolution
对相同电极几何与边界电势差 $V$，
静电位满足 $\nabla\cdot(\varepsilon\varepsilon_0
\nabla\varphi)=0$，
稳态导电满足
$\nabla\cdot(\rho^{-1}\nabla\varphi)=0$。
因为 $\rho,\varepsilon$ 均为常数，
两问题的归一化电势分布相同。
若电极在此场中通电流 $I$，
则自由电荷量 $Q$ 与电流的面积通量
分别为
\[
Q=\varepsilon\varepsilon_0\oint\bm E\cdot\mathrm d\bm S,
\qquad
I=\frac1\rho\oint\bm E\cdot\mathrm d\bm S.
\]
故 $Q/I=\rho\varepsilon\varepsilon_0$。
用 $Q=CV$、$V=IR$，得到
\[
\boxed{RC=\rho\varepsilon\varepsilon_0}.
\]
\excheck 结果与电极形状无关；
仅适用于相同空间中\emph{均匀}的
电阻率及介电常数。

\electrocurrentproblem{162}{114--115}
电阻率为 $\rho$ 的导体与相对介电常数
$\varepsilon$ 的电介质相接。
在界面某点电介质侧
位移矢量模为 $D$，指向远离导体方向，
与界面法线夹角 $\alpha$。
求导体表面自由电荷面密度
及导体内部邻近界面的稳态电流密度模。
\exsolution
从导体指向电介质取法线 $\bm n$。
法向位移的跃变为自由面电荷：
$\sigma_f=D_n=D\cos\alpha$
（理想导体法向电场内部可忽略）。
所以
\[
\boxed{\sigma_f=D\cos\alpha}.
\]
电场切向分量跨介质界面连续，
介质侧
$E_t=D_t/(\varepsilon\varepsilon_0)
=D\sin\alpha/(\varepsilon\varepsilon_0)$。
在导电体中欧姆定律
$j_t=E_t/\rho$，
且稳态电流不能穿过绝缘介质界面，
故仅有切向电流：
\[
\boxed{j=\frac{D\sin\alpha}
{\rho\varepsilon\varepsilon_0}}.
\]
\excheck 边界取 $D_n$ 决定静电面电荷，
取 $E_t$ 决定导体内沿表面的电流；
两种法向与切向条件不能互换。

\electrocurrentproblem{163}{115}
平板电极面积
$S=230\,\mathrm{cm^2}$，间距
$d=2.0\,\mathrm{mm}$；
其间弱导电介质电导率沿垂直电极方向
由 $\sigma_1=1.0\,\mathrm{pS/m}$
线性增至 $\sigma_2=2.0\,\mathrm{pS/m}$。
两极板间电压 $V=300\,\mathrm V$，
求稳态电流。
\exsolution
稳态时电流密度
$j=I/S$ 与位置无关。
局部欧姆定律
$E(x)=j/\sigma(x)$。
总电压
\[
V=j\int_0^d\frac{\mathrm dx}{\sigma_1+
(\sigma_2-\sigma_1)x/d}
=\frac{jd}{\sigma_2-\sigma_1}
\ln\frac{\sigma_2}{\sigma_1}.
\]
所以
\[
\boxed{I=\frac{VS(\sigma_2-\sigma_1)}
{d\ln(\sigma_2/\sigma_1)}}.
\]
代入 $S=0.0230\,\mathrm{m^2}$：
\[
\boxed{I\simeq5.0\times10^{-9}\,\mathrm A
=5.0\,\mathrm{nA}}.
\]
\excheck $\sigma$ 代表电导率而非面电荷密度，
其单位为 $\mathrm{S/m}$；
$\mathrm{pS/m}=10^{-12}\,\mathrm{S/m}$。

\electrocurrentproblem{164}{115}
证明两种均匀导电介质交界处稳态电流线的折射律：
若电导率分别为 $\sigma_1,\sigma_2$，
电流线与界面法线角为 $\alpha_1,\alpha_2$，
则
$\tan\alpha_2/\tan\alpha_1=\sigma_2/\sigma_1$。
\exsolution
在无电荷随时间积累的稳态边界上，
电流法向分量连续：
$j_{1n}=j_{2n}$。
由静电场切向连续：
$E_{1t}=E_{2t}$。
欧姆定律 $j_t=\sigma E_t$ 给
$j_{2t}/j_{1t}=\sigma_2/\sigma_1$。
由于
$\tan\alpha_i=j_{it}/j_{in}$，
故
\[
\frac{\tan\alpha_2}{\tan\alpha_1}
=\frac{j_{2t}}{j_{1t}}
\frac{j_{1n}}{j_{2n}}
=\boxed{\frac{\sigma_2}{\sigma_1}}.
\]
\excheck 这一折射律作用于\emph{电流线}，
不是光学斯涅尔定律；
其源自电流连续与电场切向连续。

\electrocurrentproblem{165}{115}
两段截面积相同为 $S$ 的长圆柱导体
首尾相接，电阻率分别为 $\rho_1,\rho_2$。
恒定电流 $I$ 从导体1流入导体2。
求界面上积累的总自由电荷。
\exsolution
电流密度沿圆柱轴为
$j=I/S$。
两段轴向电场分别为
$E_1=\rho_1I/S$、
$E_2=\rho_2I/S$。
交界处电场法向跃变
$E_2-E_1=\sigma_f/\varepsilon_0$，
所以
\[
\sigma_f=\frac{\varepsilon_0 I}{S}
(\rho_2-\rho_1).
\]
乘以界面面积：
\[
\boxed{q_{\rm interface}=
\varepsilon_0 I(\rho_2-\rho_1)}.
\]
\excheck 即使直流达到稳态，
不同电阻率材料的交界仍可带少量
静止表面自由电荷，以维持电场跳变。

\electrocurrentproblem{166}{115}
两层均匀介质沿电流方向串联，
厚度 $d_1,d_2$，相对介电常数
$\varepsilon_1,\varepsilon_2$，
电阻率 $\rho_1,\rho_2$。
施加直流电压 $V$，电场从1层指向2层，
求界面\emph{外来自由电荷}面密度及其为零的条件。
\exsolution
直流稳态时法向电流密度相同，
$j=E_1/\rho_1=E_2/\rho_2$，
故 $E_i=j\rho_i$。
电压
$V=j(\rho_1d_1+\rho_2d_2)$，
于是
\[
j=\frac{V}{\rho_1d_1+\rho_2d_2}.
\]
界面自由电荷密度由
$D_{2n}-D_{1n}$ 得
\[
\sigma_f=\varepsilon_0
(\varepsilon_2E_2-\varepsilon_1E_1)
=\varepsilon_0j
(\varepsilon_2\rho_2-\varepsilon_1\rho_1).
\]
所以
\[
\boxed{\sigma_f=
\frac{\varepsilon_0V
(\varepsilon_2\rho_2-\varepsilon_1\rho_1)}
{\rho_1d_1+\rho_2d_2}},
\qquad
\boxed{\sigma_f=0\iff
\varepsilon_1\rho_1=\varepsilon_2\rho_2}.
\]
\excheck 零面电荷条件不是
$\rho_1=\rho_2$ 或 $\varepsilon_1=\varepsilon_2$
单独成立，而是两个乘积相等。

\electrocurrentproblem{167}{115}
平板电极之间的弱导电介质非均匀，
靠电极1处相对介电常数和电阻率为
$\varepsilon_1,\rho_1$，
靠电极2处为 $\varepsilon_2,\rho_2$。
稳态电流 $I$ 从1流向2。
求介质\emph{体内}的总外来自由电荷。
\exsolution
取轴向坐标 $x$ 从电极1指向2，
稳态电流密度 $j=I/S$。
局部欧姆定律给 $E(x)=j\rho(x)$，
位移为
$D(x)=\varepsilon_0\varepsilon(x)j\rho(x)$。
自由体电荷满足
$\rho_f(x)=\mathrm dD/\mathrm dx$。
对介质全厚度积分：
\[
Q_f=S\int_0^d\rho_f(x)\,\mathrm dx
=S[D(d)-D(0)].
\]
因此
\[
\boxed{Q_f=
\varepsilon_0I
(\varepsilon_2\rho_2-\varepsilon_1\rho_1)}.
\]
\excheck 自由电荷总量由两边界的
$\varepsilon\rho$ 值决定，
与内部如何连续变化无关；
这里不包括两极板表面自由电荷。

\electrocurrentproblem{168}{115}
平板电极间距 $d$，
弱导电介质的电阻率沿法向线性变化，
最大与最小值之比为 $\eta>1$。
施加电压 $V$。
求间隙体自由电荷密度。
\exsolution
先按\emph{相对介电常数均匀且取1}
（与本版书后答案一致）计算。
令 $\rho(x)=\rho_{\min}
[1+(\eta-1)x/d]$，
稳态 $j$ 恒定，
$E(x)=j\rho(x)$。
电压
\[
V=j\int_0^d\rho(x)\,\mathrm dx
=\frac{jd\rho_{\min}(1+\eta)}2,
\]
故
$j=2V/[d\rho_{\min}(1+\eta)]$。
由一维高斯定律
$\rho_f=\varepsilon_0\mathrm dE/\mathrm dx$，
\[
\boxed{\rho_f=
\frac{2\varepsilon_0V(\eta-1)}
{d^2(\eta+1)}}.
\]
如果实际材料相对介电常数为\emph{均匀的}
$\varepsilon_r\ne1$，
则严格应为
$\rho_f=2\varepsilon_0\varepsilon_rV
(\eta-1)/[d^2(\eta+1)]$。
\excheck \textbf{底本条件不完整：}
原版英文最后仅写
“$\varepsilon$ is assumed to be everywhere.”
未说明其常数或取值。
书后公式不带 $\varepsilon$，相当于取
$\varepsilon_r=1$；译稿保留一般修正，
不暗称题面明确指定了1。

\electrocurrentproblem{169}{115}
长圆柱导体截面积为 $S$，
电阻率沿半径变化为
$\rho(r)=\alpha/r^2$，
$\alpha>0$。求 (a) 沿轴单位长度电阻参数；
(b) 流过总电流 $I$ 所需的轴向电场模。
\exsolution
假设轴向电场 $E$ 在同一截面各处相同，
电流密度
$j_z(r)=E/\rho(r)=Er^2/\alpha$。
若截面半径为 $R$，
\[
I=\int_Sj_z\,\mathrm dA
=\frac{2\pi E}{\alpha}\int_0^Rr^3\,\mathrm dr
=\frac{\pi ER^4}{2\alpha}.
\]
由 $S=\pi R^2$，
$I=ES^2/(2\pi\alpha)$。
因此
\[
\boxed{\frac{R_{\rm wire}}{L}
=\frac E I=\frac{2\pi\alpha}{S^2}},
\qquad
\boxed{E=\frac{2\pi\alpha I}{S^2}}.
\]
\excheck 这里 $R_{\rm wire}/L$ 为
$\Omega/\mathrm m$，
不同于第159题二维双线体系
的 $\Omega\cdot\mathrm m$ 长度参数。

\electrocurrentproblem{170}{116}
电容 $C=400\,\mathrm{pF}$
通过电阻 $R=650\,\Omega$ 接入
恒压电源 $V_0$，初始未充电。
问电容电压何时达到 $V=0.90V_0$。
\exsolution
串联 RC 充电时
$C\dot V_c=(V_0-V_c)/R$，
初值 $V_c(0)=0$。
积分：
\[
V_c(t)=V_0(1-e^{-t/(RC)}).
\]
令 $V_c/V_0=0.90$，
得 $e^{-t/(RC)}=0.10$，
故
\[
\boxed{t=-RC\ln(1-0.90)=RC\ln10}.
\]
代入
$RC=(650)(400\times10^{-12})\,\mathrm s
=2.60\times10^{-7}\,\mathrm s$，
\[
\boxed{t\simeq0.60\,\mathrm{\mu s}}.
\]
\excheck 电压趋近 $V_0$ 是渐近过程，
到达 $0.90V_0$ 需要有限时间。

\electrocurrentproblem{171}{116}
相对介电常数 $\varepsilon=2.1$ 的
电容器填充材料漏电，
电荷在 $\tau=3.0\,\mathrm{min}$
内减少到原来的一半。
假设只通过介质漏电，
求介质电阻率 $\rho$。
\exsolution
对于相同几何的电容和漏电通道，
$RC=\rho\varepsilon\varepsilon_0$。
自由放电
$Q(t)=Q_0e^{-t/(RC)}$。
在 $t=\tau$ 时 $Q/Q_0=1/2$，
所以
\[
RC=\frac{\tau}{\ln2},
\qquad
\boxed{\rho=
\frac{\tau}{\varepsilon\varepsilon_0\ln2}}.
\]
用 $\tau=180\,\mathrm s$、
$\varepsilon=2.1$，
\[
\boxed{\rho\simeq1.4\times10^{13}\,\Omega\cdot\mathrm m}.
\]
\excheck \textbf{书末数量级冲突：}
本版书末若印
$1.4\times10^{12}\,\Omega\cdot\mathrm m$，
则按题给180s、2.1和所印公式
少一个10倍；译稿按基本常数复算保留
$1.4\times10^{13}$。

\electrocurrentproblem{172}{116}
电动势恒为 $\mathcal E$ 的理想电源、
电阻 $R$ 与电容 $C$ 串联。
原电路已稳定充电，
在 $t=0$ 瞬间把电容减小为原来的
$1/\eta$，$\eta>1$。
求随后电流 $I(t)$。
\exsolution
原来稳态电容电荷
$Q(0^-)=C\mathcal E$。
电容突变瞬间电荷不及流走，
$Q(0^+)=C\mathcal E$，
但新的电容为 $C'=C/\eta$，
故电压立即跃变为
$U_C(0^+)=Q/C'=\eta\mathcal E$。
随后 RC 回路满足
$RC'\,\dot U_C=\mathcal E-U_C$，
因此
\[
U_C(t)=\mathcal E+
(\eta-1)\mathcal E\,e^{-t/(RC')}
=\mathcal E+(\eta-1)\mathcal E
e^{-\eta t/(RC)}.
\]
按从电源正端向电容充电为正，
电流 $I=(\mathcal E-U_C)/R$：
\[
\boxed{I(t)=
-\frac{(\eta-1)\mathcal E}{R}
e^{-\eta t/(RC)}}.
\]
\excheck 负号表示电容因突然减小而
\emph{倒流放电}；若以放电方向为正，
电流模为 $(\eta-1)\mathcal E
e^{-\eta t/(RC)}/R$。

\electrocurrentproblem{173}{116}
一个电动势 $\mathcal E=6.0\,\mathrm V$
的电池、安培表与伏特表串联。
随后给伏特表并联一只电阻，
伏特表读数降为原来的
$1/\eta$（$\eta=2.0$），
安培表读数升为原来的 $\eta$ 倍。
求并联电阻后伏特表读数。
\exsolution
设并联前的伏特表电压为 $U_0$，
安培表串联部分上的总电压降为 $W_0$，
电源电动势满足
$\mathcal E=U_0+W_0$。
并联后电压为 $U_0/\eta$，
串联电流增为 $\eta$ 倍，
故安培表与电源内阻合计电压降
变为 $\eta W_0$：
\[
\mathcal E=\frac{U_0}{\eta}+\eta W_0.
\]
用第一式消去 $W_0=\mathcal E-U_0$，
得
$U_0=\eta\mathcal E/(\eta+1)$，
所以新读数
\[
\boxed{U_{\rm new}=\frac{\mathcal E}{\eta+1}
=\frac6{3}\,\mathrm V=2.0\,\mathrm V}.
\]
\excheck 并不需要单独知道安培表或伏特表的
内阻值，读数倍率已足以确定电压。

\electrocurrentproblem{174}{116}
图3.39闭合回路，
$R_1=10\,\Omega$、
$R_2=20\,\Omega$、
$\mathcal E_1=5.0\,\mathrm V$、
$\mathcal E_2=2.0\,\mathrm V$，
电源内阻忽略。
求图中点1与点2的电势差。
\exsolution
图中两电源沿回路反向，
设沿 $\mathcal E_1$ 有效驱动方向为正，
回路电流
\[
I=\frac{\mathcal E_1-\mathcal E_2}{R_1+R_2}
=\frac{5-2}{10+20}=0.10\,\mathrm A.
\]
从点1沿图中相应上支路到点2，
需计电动势 $\mathcal E_1$
和电阻 $R_1$ 的代数电压变化。
由 Kirchhoff 第二定律：
\[
\varphi_1-\varphi_2
=IR_1-\mathcal E_1
=\frac{(\mathcal E_1-\mathcal E_2)R_1}
{R_1+R_2}-\mathcal E_1.
\]
所以
\[
\boxed{\varphi_1-\varphi_2=-4.0\,\mathrm V}.
\]
\excheck 负号说明点1电势比点2低4V；
电路中电源正负极必须依图3.39使用。

\electrocurrentproblem{175}{116}
两个电动势相同为 $\mathcal E$ 的电池
串联，内阻分别为 $R_1<R_2$，
外接电阻 $R$。
求使其中某一个电池的\emph{端电压}
恰为零时的外阻以及是哪一个电池。
\exsolution
两电池同向串联，
回路电流为
\[
I=\frac{2\mathcal E}{R+R_1+R_2}.
\]
供电时第 $i$ 只电池的端电压
$U_i=\mathcal E-IR_i$。
零端电压要求
$I=\mathcal E/R_i$。
由于 $R$ 必须非负，
在 $R_2>R_1$ 情况下可由
内阻较大的第2只电池达到：
\[
\frac{2\mathcal E}{R+R_1+R_2}
=\frac{\mathcal E}{R_2}
\quad\Longrightarrow\quad
\boxed{R=R_2-R_1}.
\]
\excheck 只有内阻较大的电源
$R_2$ 可在非负外阻条件下满足端电压为零；
这不表示该电池的电动势消失。
